\chapter{Spectrograms, multitapers, and periodic-versus-aperiodic spectra}\label{ch:spectral}
\begin{readergoal}Distinguish a spectral estimate from the generating spectrum and avoid inferring oscillations from arbitrary band power.\end{readergoal}
\section{Why spectral estimates fluctuate}
A finite realization contains sampling variability and boundary effects. The periodogram's frequency resolution changes with recording duration, and its variance does not simply vanish by taking a long record without appropriate averaging. Windowing reduces leakage at the price of altered resolution. Welch estimates average periodograms across segments; the segments may overlap, but this does not make them independent replicates.

\section{Spectrograms and the time-frequency trade-off}
The short-time Fourier transform is
\[X(\tau,f)=\int_{-\infty}^{\infty}x(t)w(t-\tau)e^{-i2\pi ft}\,dt,\]
where $w$ is a localized window. Narrow time windows localize fast transients while broadening spectral resolution. Long windows sharpen frequency information but smear event timing. A spectrogram $|X(\tau,f)|^2$ is not an instantaneous power measurement at a mathematical point: it summarizes a neighborhood determined by the window.

\section{Multitaper estimation}
The multitaper method forms spectral estimates using several mutually orthogonal tapers with controlled spectral concentration. Discrete prolate spheroidal sequences are commonly selected for a time-bandwidth product $NW$; the number of useful tapers is often on the order of $2NW-1$. The smoothing bandwidth and record length must be reported. Multitaper methods are central to the Chronux approach \cite{Bokil2010}.

\section{Spectral peaks over the background}
Brain spectra often contain aperiodic structure approximated locally by a $1/f^\beta$ law, along with putative narrowband components. A rise in band power may result from an oscillatory peak, a change in the aperiodic offset, a change in the exponent, or a frequency shift. Models such as FOOOF/specparam parameterize periodic and aperiodic components; IRASA uses irregular resampling to estimate the aperiodic part. These methods have distinct failure modes, especially when peaks are wide or overlap \cite{Donoghue2020,Gerster2022}.

\begin{workedexample}[Two indistinguishable band powers]
Imagine condition A has a 10-Hz Gaussian peak over a low background while condition B has no distinct peak but a sufficiently elevated broadband background. The integrals of their spectral power from 8 to 12 Hz may be equal. The same band-power summary corresponds to different generative accounts; it is therefore not a sufficient statistic for the existence of an oscillator.
\end{workedexample}
% --- expanded edition (October 2026) ---
\section{The periodogram is inconsistent}
For a stationary process sampled at $f_s$ with $N$ samples, the periodogram is
\[
\hat S(f_k)=\frac{1}{Nf_s}\Bigl|\sum_{n=0}^{N-1}x[n]e^{-2\pi ikn/N}\Bigr|^2.
\]
\begin{proposition}
If $x[n]$ is Gaussian white noise with variance $\sigma^2$, then for $0<k<N/2$ the values $\hat S(f_k)$ are independent, each distributed as $\frac{\sigma^2}{f_s}\cdot\frac{\chi^2_2}{2}$, i.e.\ exponential with mean equal to the true spectral density $\sigma^2/f_s$.
\end{proposition}
\begin{proof}[Sketch]
The real and imaginary parts of $X[k]$ are linear combinations of independent Gaussians with orthogonal coefficient vectors $\cos(2\pi kn/N)$ and $\sin(2\pi kn/N)$, each of squared norm $N/2$. They are therefore independent $\mathcal N(0,N\sigma^2/2)$ variables, and $|X[k]|^2$ is $N\sigma^2/2$ times a $\chi^2_2$ variable. Distinct bins use orthogonal coefficient vectors and are independent.
\end{proof}
The standard deviation of an exponential variable equals its mean. Doubling the record length doubles the number of bins but leaves each bin's relative error at $100\%$. This is \emph{inconsistency}: more data give finer resolution, not lower variance. Variance falls only by averaging, across segments (Welch), tapers (multitaper), neighboring frequencies (smoothing), or trials.

\section{Bias, variance, and degrees of freedom}
Averaging $K$ approximately independent periodogram-like estimates gives an estimate distributed approximately as $S(f)\chi^2_{2K}/(2K)$. Its relative standard deviation is $1/\sqrt K$, and a $95\%$ confidence interval for $S(f)$ is
\[
\Bigl[\frac{2K\hat S(f)}{\chi^2_{2K,0.975}},\ \frac{2K\hat S(f)}{\chi^2_{2K,0.025}}\Bigr],
\]
which on a logarithmic (dB) scale has a width independent of $S(f)$. This is one reason to plot spectra in decibels. Overlapping Welch segments are correlated; with $50\%$ overlap and a Hann window the effective number of independent segments is somewhat less than the nominal number.

\section{Slepian tapers}
Among all sequences of length $N$ with unit energy, which one concentrates the most of its spectral energy inside the band $[-W,W]$? The answer is the first discrete prolate spheroidal (Slepian) sequence, the leading eigenvector of the $N\times N$ matrix
\[
A_{mn}=\frac{\sin(2\pi W(m-n))}{\pi(m-n)}\quad(m\ne n),\qquad A_{nn}=2W,
\]
with $W$ in cycles per sample. The eigenvalues $\lambda_0>\lambda_1>\cdots$ are the fractions of energy inside the band. Approximately $2NW$ of them are close to one and the rest fall rapidly to zero; the corresponding eigenvectors are mutually orthogonal. The multitaper estimate averages the $K\approx2NW-1$ well-concentrated eigenspectra, achieving $K$ degrees of freedom (in pairs) with leakage controlled at resolution $2W$. A multitaper analysis therefore has two parameters with physical meaning: the window length $T=N/f_s$ and the half-bandwidth $Wf_s$ in hertz. Their product fixes $K$.

\section{The aperiodic component}
Write a parametric spectral model as
\[
\log_{10}S(f)=b-\log_{10}(k+f^\chi)+\sum_jG_j(f),
\]
with offset $b$, optional ``knee'' $k$, exponent $\chi$, and Gaussian peaks $G_j$ in log-power. When $k=0$, the aperiodic part is a straight line of slope $-\chi$ in log--log coordinates. Two consequences follow. First, a change in $\chi$ rotates the spectrum and changes band power in opposite directions at low and high frequencies, without any oscillation changing. Second, the ratio of two band powers (e.g.\ theta/delta) depends on $\chi$ even when no peak exists. The exponent has been related to the balance of excitation and inhibition in models, but that link is model-dependent and not a theorem about tissue.

\begin{workedexample}[Resolution, variance, and record length]
A 4-s LFP segment sampled at 1000\,Hz has DFT resolution $0.25$\,Hz. With multitaper half-bandwidth $W=1$\,Hz, $NW=4$ and $K=7$ tapers give $14$ degrees of freedom; the 95\% confidence interval spans roughly $-2.7$ to $+4.0$\,dB around the estimate. Doubling $W$ to $2$\,Hz yields $K=15$ and an interval of about $-2.0$ to $+2.5$\,dB, at the cost of blurring peaks narrower than $4$\,Hz. Neither choice is ``correct''; the report must state it.
\end{workedexample}
\begin{exercises}
\item Explain the meaning of time-bandwidth product in multitaper estimation.
\item Simulate a signal whose $1/f$ slope changes but whose oscillatory peak remains fixed.
\item Explain why a spectral peak detector needs a null or baseline definition.
\item Describe two ways spectrogram edge artifacts can mislead interpretation.
\item Show that the periodogram of white noise has the same relative variance for every record length.
\item Compute the 95\% confidence interval in dB for $K=1$ and $K=10$ independent averaged periodograms.
\item For $T=2$\,s and half-bandwidth $3$\,Hz, how many Slepian tapers are well concentrated?
\item Show that a change in aperiodic exponent with fixed value at $f=20$\,Hz increases 4--8\,Hz power and decreases 30--80\,Hz power.
\end{exercises}
